\documentclass[11pt]{article}
\usepackage{problem_set_style}
\usepackage{kotex}
\usepackage{booktabs}
\usepackage{array}
\begin{document}
\heading{8}{샘플}{}{Spring 2027}{중간 프로젝트 샘플 (Abstract · 보고서 · 동료평가)}

\noindent{\small 아래는 \textbf{가상의 팀}이 쓴 샘플입니다. 수치는 형식을 보이기 위한 예시이며 실제 실험 결과가 아닙니다. 형식과 분량의 기준으로만 참고하세요.}

\section*{① Abstract 샘플 (3주차 제출, 1쪽)}
\noindent\textbf{제목}: 신용카드 연체 예측에서 정규화 로지스틱회귀와 GBDT의 비교 \\
\textbf{팀}: 김OO, 이OO, 박OO

\vspace{4pt}
\noindent\textbf{문제 정의.} 행 하나는 \textbf{카드 고객 한 명}이고, 라벨은 다음 달 연체 여부(1=연체)이다. 연체 비율이 약 22\%로 한쪽에 치우쳐 있어 정확도 대신 \textbf{PR-AUC}를 주 지표로 쓴다.

\noindent\textbf{데이터.} UCI ``Default of Credit Card Clients'' (30{,}000행, 특성 23개). 같은 고객이 한 행씩만 있어 그룹 누수 위험은 없고, 시간 정보가 없어 무작위 stratified 분할(60/20/20)을 쓴다.

\noindent\textbf{베이스라인.} (a) 항상 ``연체 아님''을 내놓는 모델(정확도 78\%, PR-AUC 0.22), (b) 정규화 없는 로지스틱회귀.

\noindent\textbf{계획.} L2 정규화 로지스틱회귀와 GBDT를 비교한다. 정규화 세기와 트리 개수는 validation 곡선으로 고르고, test는 마지막에 한 번만 측정한다. 수식적 정당화로는 정규화에 따른 편향-분산 변화와 GBDT의 SHAP 분해 한 건을 쓴다.

\noindent\textbf{예상되는 위험.} 청구 금액 특성들이 서로 강하게 연관되어 로지스틱회귀 계수가 불안정할 수 있다. 정규화로 계수가 어떻게 변하는지 함께 기록한다.

\newpage
\section*{② 보고서 샘플 (6주차 제출, 축약본)}
\noindent\textbf{1. 문제 정의.} 고객 한 명을 행으로, 다음 달 연체 여부를 라벨로 한다. 주 지표는 PR-AUC이다(연체 22\%의 불균형 때문). 베이스라인은 항상 ``연체 아님'' (PR-AUC 0.22).

\noindent\textbf{2. 데이터와 전처리.} 30{,}000행을 60/20/20으로 stratified 분할했다. 표준화의 평균과 표준편차는 \textbf{train에서만} 구해 validation과 test에 적용했다. 범주형(학력, 결혼 상태)은 원-핫 인코딩.

\noindent\textbf{3. 모델 선택과 이유.} 로지스틱회귀는 해석이 쉽고, 특성이 많이 연관되어 있어 L2 정규화 효과를 보이기 좋다. GBDT는 표 데이터에서 강한 비교 대상이다. 거리 기반이 아니므로 GBDT에는 표준화를 쓰지 않았다.

\begin{center}\small
\begin{tabular}{@{}lcc@{}}\toprule
모델 (validation) & PR-AUC & 비고 \\\midrule
베이스라인(항상 0) & 0.22 & \\
로지스틱회귀, 정규화 없음 & 0.46 & \\
로지스틱회귀, L2 ($C=0.1$) & 0.48 & validation으로 선택 \\
GBDT (트리 150개) & 0.55 & validation으로 선택 \\\bottomrule
\end{tabular}
\end{center}

\noindent\textbf{4. 결과.} 최종 test는 \textbf{한 번만} 측정했다: GBDT PR-AUC 0.54, 로지스틱회귀(L2) 0.47. 학습 곡선(그림 1)에서 트리 개수가 150을 넘으면 train은 계속 오르지만 validation은 멈춰 150에서 멈췄다.

\noindent\textbf{5. 수식적 정당화.} L2 정규화가 계수의 크기를 줄여 분산을 낮춘 결과 validation PR-AUC가 0.46에서 0.48로 올랐다(편향은 약간 증가). SHAP으로 한 고객의 예측을 분해하면 최근 두 달의 연체 여부가 가장 크게 기여했다.

\noindent\textbf{6. 잘 안 된 부분.} 불균형을 줄이려고 연체 쪽 가중치를 키웠더니 PR-AUC가 오히려 0.02 내려갔다. 가중치가 확률을 왜곡해 순위 지표에는 도움이 안 된다고 진단했다.

\newpage
\section*{③ 동료평가 샘플 (8주차 제출)}
\noindent\textbf{평가 대상}: 다른 팀의 ``당뇨 재입원 예측'' 보고서 \quad \textbf{작성자}: 이OO

\begin{center}\small
\begin{tabular}{@{}p{2.6cm}p{9.8cm}@{}}\toprule
항목 & 평가 (좋았던 점 / 아쉬운 점) \\\midrule
문제 정의 & 환자 한 명이 아니라 \textbf{입원 한 건}이 행인데, 한 환자의 여러 입원이 train과 test에 섞였는지 설명이 없다. 라벨과 지표는 명확하다. \\
방법론 적절성 & 클래스 비율(11\%)을 보고 PR-AUC를 고른 것은 적절하다. 다만 거리 기반 모델에 표준화를 했는지 적혀 있지 않다. \\
수식적 정당화 & SHAP 분해로 ``왜 이 환자가 위험한가''를 보였고 설명이 실제 숫자와 맞는다. \\
결과의 정직성 & ``test를 한 번만 썼다''고 명시했고 실패한 시도도 적었다. 좋다. \\\bottomrule
\end{tabular}
\end{center}

\noindent\textbf{반박 가능한 질문.} 같은 환자의 입원이 train과 test에 함께 들어갔다면 test PR-AUC 0.31은 부풀려진 숫자가 아닐까요? 환자 단위로 나눠서 다시 측정한 결과가 있나요?

\end{document}
